% ==============================================================================
% Slide 11: iForest Geometry & Subspace Partitioning
% ==============================================================================
\begin{frame}{iForest Geometry \& Subspace Partitioning}{Recursive Hyperplane Cuts Isolate Anomalies at Shallow Tree Depths}
    \vspace{-0.1cm}
    \begin{columns}[T,onlytextwidth]
        % Column 1: 2D Feature Space Cuts Diagram
        \begin{column}{0.485\textwidth}
            \begin{tcolorbox}[
                enhanced,
                colback=cardbg,
                colframe=mintborder,
                boxrule=0.7pt,
                arc=3pt,
                height=5.45cm,
                left=4pt, right=4pt, top=4pt, bottom=4pt
            ]
                \centering
                {\fontsize{7.5pt}{9pt}\selectfont\bfseries\color{deepemerald}2D Feature Space Random Hyperplanes}\par
                \vspace{4pt}
                \begin{tikzpicture}[scale=0.88, every node/.style={transform shape}]
                    % 2D Feature Space Boundary Box
                    \draw[draw=mintborder, fill=cardinner, line width=0.8pt, rounded corners=2pt] (-2.9, 0.4) rectangle (2.9, 3.2);
                    \node[font=\fontsize{5.5pt}{6.5pt}\selectfont\bfseries\color{slatygray}] at (0, 3.0) {Ratio Subspace: Operating Margin vs.\ OCF / Assets};

                    % Normal peer cluster
                    \foreach \x/\y in {-0.2/1.6, 0.2/1.8, -0.4/1.9, 0.5/1.5, 0.1/1.4, -0.1/2.0, 0.4/2.1, -0.5/1.7, 0.0/1.7, 0.3/1.9} {
                        \node[circle, fill=forestgreen, inner sep=1.4pt] at (\x, \y) {};
                    }
                    \node[below, font=\fontsize{5pt}{6pt}\selectfont\color{forestgreen}] at (0.0, 1.3) {Normal Peer Cluster};

                    % Anomaly in upper left corner
                    \node[circle, fill=jade, inner sep=2.2pt] (anom) at (-2.3, 2.6) {};
                    \node[above, font=\fontsize{5.2pt}{6.2pt}\selectfont\bfseries\color{deepemerald}] at (-2.3, 2.65) {$x_{\text{outlier}}$};

                    % Random Hyperplane Cuts
                    % Cut 1: Horizontal cut at y=2.2
                    \draw[line width=0.9pt, dashed, color=deepemerald] (-2.9, 2.2) -- (2.9, 2.2)
                        node[right, font=\fontsize{4.8pt}{5.8pt}\selectfont\color{deepemerald}] {Cut 1};
                    % Cut 2: Vertical cut at x=-1.6 above Cut 1
                    \draw[line width=0.9pt, dash dot, color=jade] (-1.6, 2.2) -- (-1.6, 3.2)
                        node[above, font=\fontsize{4.8pt}{5.8pt}\selectfont\color{jade}] {Cut 2};

                    % Multiple cuts required for normal points (faint lines)
                    \draw[line width=0.5pt, dotted, color=slatygray] (-1.0, 0.4) -- (-1.0, 2.2);
                    \draw[line width=0.5pt, dotted, color=slatygray] (1.2, 0.4) -- (1.2, 2.2);
                    \draw[line width=0.5pt, dotted, color=slatygray] (-1.0, 1.2) -- (1.2, 1.2);

                    % Annotation Box
                    \node[fill=cardinner, draw=mintborder, rounded corners=2pt, inner sep=3.5pt, text width=5.2cm, align=left] at (0, -0.2) {
                        \fontsize{5.2pt}{6.5pt}\selectfont
                        \textbf{\color{deepemerald}Rapid Partitioning Principle:}
                        Because $x_{\text{outlier}}$ is isolated in low-density space, only \textbf{2 random bisections} separate it completely from the peer cohort.
                    };
                \end{tikzpicture}
                \vfill
                \badge[jade]{Subspace Hyperplane Bisection Mechanics}
            \end{tcolorbox}
        \end{column}

        % Column 2: Isolation Tree Hierarchy & Ensemble Properties
        \begin{column}{0.485\textwidth}
            \begin{tcolorbox}[
                enhanced,
                colback=cardbg,
                colframe=mintborder,
                boxrule=0.7pt,
                arc=3pt,
                height=2.65cm,
                left=4pt, right=4pt, top=4pt, bottom=4pt
            ]
                \centering
                {\fontsize{7.5pt}{9pt}\selectfont\bfseries\color{deepemerald}iTree Hierarchy: Shallow Depth Isolation}\par
                \vspace{2pt}
                \begin{tikzpicture}[scale=0.88, every node/.style={transform shape}]
                    % Root
                    \node[draw=deepemerald, fill=cardinner, rounded corners=2pt, font=\fontsize{5.2pt}{6.2pt}\selectfont\bfseries, inner sep=2.5pt, text=darkcharcoal] (root) at (0, 1.7) {Root (Cut 1)};
                    
                    % Children
                    \node[draw=jade, fill=jade!25!cardbg, rounded corners=2pt, font=\fontsize{5.2pt}{6.2pt}\selectfont\bfseries, inner sep=2.5pt, text=deepemerald, below left=0.45cm and 0.8cm of root] (c1) {Cut 2};
                    \node[draw=mintborder, fill=cardinner, rounded corners=2pt, font=\fontsize{5.2pt}{6.2pt}\selectfont, inner sep=2.5pt, text=darkcharcoal, below right=0.45cm and 0.8cm of root] (c2) {Dense Partition};

                    \draw[->, >=Stealth, line width=0.7pt, color=deepemerald] (root) -- (c1);
                    \draw[->, >=Stealth, line width=0.7pt, color=deepemerald] (root) -- (c2);

                    % Outlier leaf
                    \node[draw=jade, fill=jade!40!cardbg, rounded corners=2pt, font=\fontsize{5.2pt}{6.2pt}\selectfont\bfseries, inner sep=2.5pt, text=deepemerald, below left=0.45cm and 0.3cm of c1] (leaf1) {Isolated $x_{\text{outlier}}$};
                    \node[font=\fontsize{5.2pt}{6.2pt}\selectfont\color{slatygray}, below right=0.45cm and 0.3cm of c1] (other) {$\dots$};

                    \draw[->, >=Stealth, line width=0.7pt, color=jade] (c1) -- node[left, font=\fontsize{4.8pt}{5.8pt}\selectfont\bfseries\color{jade}] {$h = 2$} (leaf1);
                    \draw[->, >=Stealth, line width=0.7pt, color=slatygray] (c1) -- (other);

                    % Deep nodes
                    \node[font=\fontsize{5.2pt}{6.2pt}\selectfont\color{slatygray}, below=0.45cm of c2] (deep) {Deep Leaves ($h \ge 8$)};
                    \draw[->, >=Stealth, line width=0.7pt, dashed, color=slatygray] (c2) -- (deep);
                \end{tikzpicture}
            \end{tcolorbox}

            \vspace{0.15cm}

            \begin{quantcard}[height=2.65cm]{Forest Ensemble Convergence}
                \fontsize{6.3pt}{8.2pt}\selectfont
                \textbf{\color{deepemerald}Path Length Averaging Across Trees}
                \begin{equation*}
                    \mathbb{E}[h(x)] = \frac{1}{T} \sum_{t=1}^{T} h_t(x), \quad T = 100 \text{ Trees}
                \end{equation*}
                \begin{itemize}
                    \setlength{\itemsep}{1.5pt}
                    \item Averaging over 100 stochastic trees smooths out random split noise and yields monotonic convergence of path length.
                    \item \textbf{Curse of Dimensionality Immunity:} Random subsampling preserves geometric distance without requiring expensive matrix inversions.
                \end{itemize}
            \end{quantcard}
        \end{column}
    \end{columns}
\end{frame}
